\clearpage\thispagestyle{empty}
\pdfbookmark[0]{Quelle zu Vesta (8. April 1809)}{vesta1809apr-sources}
\begin{center}\textbf{Quelle zu Vesta (8. April 1809)}\end{center}
\noindent Diese KI-Übertragung ins heutige Deutsch folgt der Mitteilung vom 8. April 1809 in Gauß’ \href{https://gdz.sub.uni-goettingen.de/id/PPN236007467}{\textit{Werke}, Bd.~VI} (Göttingen, 1874), S.~309--311. Der Erstdruck erschien unter \textit{Göttingen} in den \textit{Göttingischen gelehrten Anzeigen}, Stück 55, S.~\href{https://www.digitale-sammlungen.de/en/view/bsb10538675?page=539}{537}, \href{https://www.digitale-sammlungen.de/en/view/bsb10538675?page=540}{538}, \href{https://www.digitale-sammlungen.de/en/view/bsb10538675?page=541}{539} und \href{https://www.digitale-sammlungen.de/en/view/bsb10538675?page=542}{540}.
\par\medskip
Die Mitteilung enthält die Pariser und Göttinger Beobachtungen, die erste beobachtete Opposition der Vesta, ihre vierten Bahnelemente und Angaben zur relativen Helligkeit. Der letzte Punkt der Punktreihe in der Epochenzeile 1811 steht wie in der Werkausgabe erhöht. KI-Übertragung und KI-Satz.
\clearpage
